\section{Usage}

It is \textbf{strongly recommended} to use \texttt{sync-template}, which is a CLI tool that can automatically synchronize the template updates.
Visit \url{https://github.com/XXCC-Unicorn-Press/sync-template} for more details.

\subsection{Initialize a new project}

First you need to initialize a new project with the template.
You can do this by running the following command in your terminal:
\begin{tcblisting}{listing only, sharp corners, boxrule=0.5pt, colback=white}
sync-template init gh:XXCC-Unicorn-Press/article-template <project-directory>
\end{tcblisting}
Alternatively, you can also clone the template repository your self by running the following command: 
\begin{tcblisting}{listing only, sharp corners, boxrule=0.5pt, colback=white}
  git clone --depth 1 https://github.com/XXCC-Unicorn-Press/article-template.git <project-directory>
\end{tcblisting}

\subsection{Build the project}

After initializing the project, you can build it by running \verb|latexmk| in the project directory.
You will get a PDF file named \texttt{main.pdf} in the project directory after the build is complete.
The default result is exactly this document you are reading.

\subsection{Modify the metadata}

You can modify the metadata in \texttt{preamble/meta.tex}.
See \Cref{tab:metadata} for the description of each metadata.

\begin{table}[!ht]
  \centering

  \begin{tabular}{lp{0.5\textwidth}}
    \hline\hline
    \textbf{Metadata} & \textbf{Description} \\
    \hline
    \verb|\doctitle| & The title displayed on the title page. You can use text formats or math equations. \\
    \verb|\doctitlestring| & The title written in the PDF metadata. This should be a pure string. \\
    \verb|\docauthor| & The author. This should be a pure string \\
    \hline\hline
  \end{tabular}

  \caption{Metadata in \texttt{preamble/meta.tex}.}
  \label{tab:metadata}
\end{table}

\subsection{Write the content}

The content of the paper is organized in the \texttt{content} directory.
It is a good habit to create a new file per section and name the file after the section name.

You should \verb|\input| the content files in \texttt{main.tex} in the order they appear in the paper.

\subsection{Insert a figure}

To insert a figure, first place the figure file in the \texttt{figures} directory, and then use the \verb|\includegraphics| command to include the figure in the content file.

If you are using Inkscape to create your figures, you can place the source files in \texttt{figures/inkscape} and run \verb|./figures/inkscape/export.sh| to export all the \verb|.svg| files to \verb|.png| files, which can be included in the document directly.

See \Cref{sec:convenient-figure-management} for more details on figure management.

\subsection{Synchronize template updates}

Sometimes the template may be updated, but your local copy will not be automatically synchronized.

If you are using \texttt{sync-template}, you can simply run \verb|sync-template sync| in the project directory to synchronize the updates.
Refer to \url{https://github.com/XXCC-Unicorn-Press/sync-template} for more details.

If you are not using \texttt{sync-template}, you can manually examine the updates in the template repository and apply the necessary changes to your local copy.
